% Generated from saved Week 3 outputs; no model calls.
\begin{table*}[t]
\centering
\small
\setlength{\tabcolsep}{6pt}
\renewcommand{\arraystretch}{1.14}
\begin{tabular}{lrr}
\toprule
Mechanical audit & GPT-5-mini & Qwen3.6-27B \\
\midrule
Supplied evidence spans & 300 & 167 \\
Located in named input & 167/300 (55.7\%) & 129/167 (77.2\%) \\
Unmatched spans & 133 & 38 \\
Cases with audit flags & 32/37 (86.5\%) & 21/37 (56.8\%) \\
\bottomrule
\end{tabular}
\caption{Citation-location checks on the same 37 cases. Span rates pool all supplied evidence units; case rates use 37 as denominator.}
\label{tab:cgen-t05}
\vspace{2pt}
\begin{minipage}{\textwidth}
\footnotesize The existing audit lowercases text and extracts ASCII alphanumeric tokens before searching for contiguous spans of at least three tokens in the named input document. Thus located spans are token-normalized matches, not verified literal quotations. Counts include claim-ground evidence and applicant-argument support spans; duplicate spans are not deduplicated. Location does not establish that evidence supports an argument. The models supply different numbers of spans, so this is not a recall measure or proof of overall evidence quality.
\end{minipage}
\end{table*}
